\chapter{总结与展望}

\section{全文总结}

本文围绕三维点云几何感知与机器人抓取规划领域长期存在的感知精度受限与接触模式相对单一两类关键问题，构建并验证了一条由几何特征检测通向复杂接触抓取生成的完整技术路径。全文首先在问题层面阐明了边缘与角点信息在三维结构表达和抓取决策中的核心作用，继而在方法层面提出兼具学习能力、可解释性与工程可实施性的解决方案，最终依托数据集评测、仿真评估与真机实验形成系统化证据链，对方法有效性进行了全面检验。

在几何特征检测方面，本文提出了 NECDNet 框架，通过多尺度邻域建模与显式解析特征融合，实现了面点、边缘点与角点的统一检测。该方法并未单纯依赖黑箱式端到端表示学习，而是在局部几何建模阶段引入曲率、投影距离均值、中心偏离比、法向角度差、法向分布方差与法向投影距离均值等具有明确物理意义的统计量，并在 $k\in\{8,16,32,64\}$ 的多尺度空间上组织为结构化输入特征矩阵，从而兼顾了模型判别能力与结果可解释性。网络结构采用分层分类思想，先区分面点与非面点，再细分边缘点与角点，配合 Focal Loss 与针对角点稀缺性的重采样策略，有效缓解了类别极不均衡场景下的训练偏置问题。基于 ABC-Corner 数据集的对比实验表明，NECDNet 在多项核心指标上均取得领先结果，其中使用网格法向量输入时宏平均 F1 达到 0.815、MCC 达到 0.943；在二分类边缘检测设置下，NECDNet-2c$\dagger$ 的 F1 达到 0.956。上述结果说明，本文方法不仅提升了边缘检测精度，更在传统方法普遍薄弱的角点检测任务上取得了显著收益。

在数据基础设施方面，本文围绕高质量角点标注稀缺这一制约学习方法性能上限的关键问题，构建了 ABC-Corner 数据集。该数据集以 ABC 的 CAD 模型为基础，形成了规则驱动自动粗标注与人工精修校正相结合的双阶段标注流程，并在样本筛选与质量控制中强调几何特征代表性和标注一致性。通过补充角点精标注，数据集将原有以面和边为主的标注体系扩展为覆盖面、边、角三类特征的完整监督体系，不仅支撑了 NECDNet 的训练与评估，也为后续研究提供了可复现、可比较的实验基准。

在抓取规划方面，本文进一步将几何检测结果转化为抓取决策约束，提出了考虑边缘与角点接触的六自由度抓取规划方法。针对传统规划器主要依赖平面接触且难以覆盖复杂几何外形的问题，本文通过类别感知体素采样缓解候选接触点在面区域的长尾聚集现象，并为面、边、角分别设计圆柱、扇形柱与圆锥检测器，实现与接触自由度一致的搜索空间建模。随后，方法按照六类接触组合构建差异化配对规则，在对跖约束、投影约束与几何一致性约束的联合作用下生成高质量候选接触对，并结合绕接触轴旋转增强、碰撞检测与软接触稳定性评估形成完整流程。相较于将软指接触简化为单点摩擦模型的经典处理，本文采用离散接触簇近似连续变形接触，能够在可接受计算代价下更准确地表征边缘与角点区域的受力特性。

在实验验证方面，本文通过第五章建立了覆盖仿真与真机场景的系统评估框架。抓取分布多样性分析显示，本文方法在代表性物体上的 STD、R 与 MAD 指标均达到最优，表明其生成的抓取解集在空间覆盖性与方向分散性方面显著优于对比方法。物理抓取实验覆盖十类不同几何特征物体，本文方法平均成功率达到 97.5\%，明显高于 GraspIt!、GPD 与 Wan's，且在无平行面或局部曲率变化剧烈的对象上仍保持稳定性能。效率评估结果进一步表明，单次规划平均耗时约 1.583 s，具备工程应用所需的时效性。综合而言，本文完成了由几何感知能力提升到抓取行为质量提升的因果闭环验证，证明了边缘与角点接触建模在复杂对象抓取中的实际价值。

\section{创新点}

本文的创新性主要体现在方法链条的完整性、关键环节的机制设计以及跨任务验证的一致性。首先，在点云几何特征检测任务中，本文提出了将显式几何统计与轻量级学习架构深度耦合的 NECDNet，将多尺度局部几何信息组织为结构化特征输入，并通过分层分类机制提升边缘与角点的可分性。该设计避免了单纯依赖大规模参数学习带来的解释性不足问题，也突破了传统阈值法在复杂场景中的泛化瓶颈，特别是在角点这一高稀疏、高混淆类别上的性能改进具有明显方法学意义。

其次，本文构建了覆盖面、边、角三类标签的 ABC-Corner 数据集，形成了可复现的数据生成与校核流程。该数据集并非仅作为训练样本集合存在，而是在研究中承担了方法开发、性能对比与消融分析的统一评测基座作用，使算法性能提升能够建立在稳定、可追溯的标注标准之上。通过这一数据层创新，本文为相关研究由经验性比较走向规范化比较提供了重要支撑。

再次，本文在抓取规划中显式引入边缘与角点接触，提出与几何特征自由度相匹配的检测器建模和配对策略，并结合软接触离散化稳定性评估，实现了从候选生成到力学可行性判定的统一建模框架。与强调面接触保守性的传统规划路线相比，本文方法在保证稳定性的同时显著扩展了抓取解空间覆盖范围，使规划器在非平行面、曲面与复杂齿状结构对象上的可用性得到实质性提升。该创新不局限于某一独立模块性能的提高，而是体现为感知、几何推理与接触力学三者协同优化后的系统级性能跃迁。

\section{研究不足与展望}
本文工作在取得较高检测精度与抓取成功率的同时，也暴露出若干值得继续深入的问题。首先，NECDNet 对法向量质量仍具有一定敏感性。实验已表明，精确网格法向量可显著提升检测上限，而在实际扫描点云中，局部最小二乘估计在尖锐区域的平滑化偏差会压缩边缘与角点间的特征间隔。这意味着后续研究需要进一步降低模型对高质量外部法向先验的依赖，可通过引入更鲁棒的法向估计模型、构建面向尖锐结构的可微法向优化模块，或直接将法向估计与特征分类联合训练，在同一目标函数下实现误差反向传播与协同收敛。

其次，本文抓取规划框架仍建立在已知完整模型的前提之上，更适用于离线规划与相对受控场景。面对遮挡严重、对象未知或仅有部分可见点云的在线作业环境，现有流程的前提条件将被削弱。后续可将本研究与增量式三维重建、点云补全和不确定性建模结合，在局部几何可观测信息不足时进行风险感知规划，并通过闭环视觉反馈动态修正候选抓取集合，从而推动方法由“完整模型驱动”向“部分观测驱动”迁移。

再次，当前方法主要面向双指平行夹爪，尽管在复杂几何对象上已展现出明显优势，但尚未充分覆盖多指灵巧手场景。多指协同接触不仅涉及更高维的接触组合空间，也涉及手内重构、接触序列切换与力分配协调等问题。未来可在本文面、边、角接触建模思想基础上，研究面向多指系统的接触图优化与协同约束求解，以支持更复杂的操作任务。

最后，本文的优化目标主要聚焦几何稳定性与执行可行性，对任务语义与功能可供性约束的考虑仍显不足。在实际工业与服务场景中，抓取行为通常受后续任务链条制约，例如工具类物体需优先保留功能端，装配类任务需满足特定姿态约束。后续研究可进一步将语义理解、任务规划与接触几何统一于同一决策框架，使抓取策略由几何可抓提升至任务可用、过程可执行且结果可交付的更高层次智能操作能力。

